% --- Archon physics formalization source begin ---
% archon:physics
% archon:covers PhyXMiniProblems/problem_phyx_mini_0275.lean
% archon:source-report reports/phyx_mini/problem_phyx_mini_0275.source.json

\chapter{Physics problem phyx\_mini\_0275}
\label{ch:PhyXMiniProblems_problem_phyx_mini_0275}

\paragraph{Problem source.}
\#\# Physical scenario

A simple harmonic oscillator consists of a \$0.50\textbackslash{} kg\$ block attached to a spring. The block slides back and forth along a straight line on a frictionless surface with equilibrium point \$x = 0\$. At \$t = 0\$ the block is at \$x = 0\$ and moving in the positive \$x\$ direction. A graph of the magnitude of the net force \$\textbackslash{}vec\{F\}\$ on the block as a function of its position is shown in figure. The vertical scale is set by \$F\_s = 75.0\textbackslash{} N\$.

\#\# Question

What is the maximum kinetic energy?

\#\# Answer choices

- A: A: 10.2 J

- B: B: 10.8 J

- C: C: 11.3 J

- D: D: 12.4 J

\#\# Figure caption (auxiliary; use the image as primary evidence)

The image depicts a graph with the x-axis labeled "x (m)" and the y-axis labeled "F (N)." The graph shows a piecewise linear function with two segments forming a V shape.

1. **Axes:**
   - The horizontal axis (x-axis) is marked with values -0.30 and 0.30 meters.
   - The vertical axis (y-axis) is labeled with "F (N)" for force in newtons.

2. **Graph Lines:**
   - There is a diagonal line starting from the point (-0.30, \textbackslash{}(F\_s\textbackslash{})), decreasing to the origin (0, 0), and continuing down to the point (0.30, -\textbackslash{}(F\_s\textbackslash{})).

3. **Dashed Lines:**
   - Horizontal dashed lines extend from the y-axis to the ends of the V shape, indicating the force values \textbackslash{}(F\_s\textbackslash{}) and -\textbackslash{}(F\_s\textbackslash{}).
   - Vertical dashed lines extend from the ends of the V shape down to the x-axis at -0.30 and 0.30.

4. **Annotations:**
   - The points where the dashed lines intersect the axes are labeled with values such as \textbackslash{}(F\_s\textbackslash{}), -\textbackslash{}(F\_s\textbackslash{}), 0.30, and -0.30.

The graph represents a force vs. displacement scenario, where \textbackslash{}(F\_s\textbackslash{}) is the maximum positive and negative force applied at corresponding displacements of -0.30 and 0.30 meters.

\#\# Dataset metadata

Wave Properties; Physical Model Grounding Reasoning, Multi-Formula Reasoning

\paragraph{Recorded answer/context.}
C: C: 11.3 J

\paragraph{Figure/image path.}
/root/proposal\_for\_physic/science-mango/phyx\_mini\_run/phyx\_data/test\_image/275.png

\paragraph{Formalization target.}
create a compiling Lean file with sorry bodies at `PhyXMiniProblems/problem\_phyx\_mini\_0275.lean`. The Lean declarations must preserve the physical quantities, dimensions or dimensional roles, figure labels, governing-law hypotheses, and final relation expressed by this problem.
Use Mathlib/Physlib names found through LeanExplore where available. If a physics API is missing, introduce faithful local abstractions rather than scalar placeholder aliases.

\begin{theorem}[Physics formalization target]
\label{thm:physics:phyx_mini_0275:target}
\lean{PhyXMiniProblems.ProblemPhyXMini0275.problem_phyx_mini_0275}
\uses{def:physics:phyx-mini-0275:phyxminiproblems-problemphyxmini0275-energyinjoules, def:physics:phyx-mini-0275:phyxminiproblems-problemphyxmini0275-blockspringoscillatorsetup, def:physics:phyx-mini-0275:phyxminiproblems-problemphyxmini0275-matchesproblemdescription, def:physics:phyx-mini-0275:phyxminiproblems-problemphyxmini0275-matchesprimaryfigure, def:physics:phyx-mini-0275:phyxminiproblems-problemphyxmini0275-hasphysicalparameters, def:physics:phyx-mini-0275:phyxminiproblems-problemphyxmini0275-satisfiesundampedsimpleharmonicoscillatorlaws, def:physics:phyx-mini-0275:phyxminiproblems-problemphyxmini0275-ismaximumkineticenergy, def:physics:phyx-mini-0275:phyxminiproblems-problemphyxmini0275-answerchoice, def:physics:phyx-mini-0275:phyxminiproblems-problemphyxmini0275-roundstonearesttenth, def:physics:phyx-mini-0275:phyxminiproblems-problemphyxmini0275-isclosestdisplayedanswer}
The assigned autoformalize agent should translate this physics problem into Lean declarations in the covered file, with theorem and lemma proof bodies written as `by sorry`.
\end{theorem}
\begin{proof}
This is an autoformalization task, not a proof task. Produce statements that can later be proved by the physics prover without weakening the physical model.
\end{proof}

% --- Archon declaration topology begin ---
\section{Declaration topology}
\begin{definition}[Mass Quantity]
  \label{def:physics:phyx-mini-0275:phyxminiproblems-problemphyxmini0275-massquantity}
  \lean{PhyXMiniProblems.ProblemPhyXMini0275.MassQuantity}
  A physical mass
\end{definition}
\begin{proof}
  This is the declared model definition of Mass Quantity.
\end{proof}
\begin{definition}[Length Quantity]
  \label{def:physics:phyx-mini-0275:phyxminiproblems-problemphyxmini0275-lengthquantity}
  \lean{PhyXMiniProblems.ProblemPhyXMini0275.LengthQuantity}
  A signed one-dimensional position or positive amplitude
\end{definition}
\begin{proof}
  This is the declared model definition of Length Quantity.
\end{proof}
\begin{definition}[Time Quantity]
  \label{def:physics:phyx-mini-0275:phyxminiproblems-problemphyxmini0275-timequantity}
  \lean{PhyXMiniProblems.ProblemPhyXMini0275.TimeQuantity}
  A physical time coordinate
\end{definition}
\begin{proof}
  This is the declared model definition of Time Quantity.
\end{proof}
\begin{definition}[Velocity Quantity]
  \label{def:physics:phyx-mini-0275:phyxminiproblems-problemphyxmini0275-velocityquantity}
  \lean{PhyXMiniProblems.ProblemPhyXMini0275.VelocityQuantity}
  A signed velocity component along the oscillator's x axis
\end{definition}
\begin{proof}
  This is the declared model definition of Velocity Quantity.
\end{proof}
\begin{definition}[Force Component Quantity]
  \label{def:physics:phyx-mini-0275:phyxminiproblems-problemphyxmini0275-forcecomponentquantity}
  \lean{PhyXMiniProblems.ProblemPhyXMini0275.ForceComponentQuantity}
  A signed force component along the oscillator's x axis
\end{definition}
\begin{proof}
  This is the declared model definition of Force Component Quantity.
\end{proof}
\begin{definition}[Spring Stiffness Quantity]
  \label{def:physics:phyx-mini-0275:phyxminiproblems-problemphyxmini0275-springstiffnessquantity}
  \lean{PhyXMiniProblems.ProblemPhyXMini0275.SpringStiffnessQuantity}
  Spring stiffness, with dimension force divided by length
\end{definition}
\begin{proof}
  This is the declared model definition of Spring Stiffness Quantity.
\end{proof}
\begin{definition}[mass In Kilograms]
  \label{def:physics:phyx-mini-0275:phyxminiproblems-problemphyxmini0275-massinkilograms}
  \lean{PhyXMiniProblems.ProblemPhyXMini0275.massInKilograms}
  \uses{def:physics:phyx-mini-0275:phyxminiproblems-problemphyxmini0275-massquantity}
  Kilogram readout of a physical mass
\end{definition}
\begin{proof}
  This is the declared model definition of mass In Kilograms.
\end{proof}
\begin{definition}[length In Meters]
  \label{def:physics:phyx-mini-0275:phyxminiproblems-problemphyxmini0275-lengthinmeters}
  \lean{PhyXMiniProblems.ProblemPhyXMini0275.lengthInMeters}
  \uses{def:physics:phyx-mini-0275:phyxminiproblems-problemphyxmini0275-lengthquantity}
  Metre readout of a signed position or amplitude
\end{definition}
\begin{proof}
  This is the declared model definition of length In Meters.
\end{proof}
\begin{definition}[time In Seconds]
  \label{def:physics:phyx-mini-0275:phyxminiproblems-problemphyxmini0275-timeinseconds}
  \lean{PhyXMiniProblems.ProblemPhyXMini0275.timeInSeconds}
  \uses{def:physics:phyx-mini-0275:phyxminiproblems-problemphyxmini0275-timequantity}
  Second readout of a physical time
\end{definition}
\begin{proof}
  This is the declared model definition of time In Seconds.
\end{proof}
\begin{definition}[velocity In Meters Per Second]
  \label{def:physics:phyx-mini-0275:phyxminiproblems-problemphyxmini0275-velocityinmeterspersecond}
  \lean{PhyXMiniProblems.ProblemPhyXMini0275.velocityInMetersPerSecond}
  \uses{def:physics:phyx-mini-0275:phyxminiproblems-problemphyxmini0275-velocityquantity}
  Metres-per-second readout of a signed velocity component
\end{definition}
\begin{proof}
  This is the declared model definition of velocity In Meters Per Second.
\end{proof}
\begin{definition}[force Component In Newtons]
  \label{def:physics:phyx-mini-0275:phyxminiproblems-problemphyxmini0275-forcecomponentinnewtons}
  \lean{PhyXMiniProblems.ProblemPhyXMini0275.forceComponentInNewtons}
  \uses{def:physics:phyx-mini-0275:phyxminiproblems-problemphyxmini0275-forcecomponentquantity}
  Newton readout of a signed force component
\end{definition}
\begin{proof}
  This is the declared model definition of force Component In Newtons.
\end{proof}
\begin{definition}[spring Stiffness In Newtons Per Meter]
  \label{def:physics:phyx-mini-0275:phyxminiproblems-problemphyxmini0275-springstiffnessinnewtonspermeter}
  \lean{PhyXMiniProblems.ProblemPhyXMini0275.springStiffnessInNewtonsPerMeter}
  \uses{def:physics:phyx-mini-0275:phyxminiproblems-problemphyxmini0275-springstiffnessquantity}
  Newton-per-metre readout of a physical spring stiffness
\end{definition}
\begin{proof}
  This is the declared model definition of spring Stiffness In Newtons Per Meter.
\end{proof}
\begin{definition}[energy In Joules]
  \label{def:physics:phyx-mini-0275:phyxminiproblems-problemphyxmini0275-energyinjoules}
  \lean{PhyXMiniProblems.ProblemPhyXMini0275.energyInJoules}
  Joule readout of Physlib's dimensionful energy quantity
\end{definition}
\begin{proof}
  This is the declared model definition of energy In Joules.
\end{proof}
\begin{definition}[Restoring Element]
  \label{def:physics:phyx-mini-0275:phyxminiproblems-problemphyxmini0275-restoringelement}
  \lean{PhyXMiniProblems.ProblemPhyXMini0275.RestoringElement}
  The component that provides the restoring force
\end{definition}
\begin{proof}
  This is the declared model definition of Restoring Element.
\end{proof}
\begin{definition}[Surface Condition]
  \label{def:physics:phyx-mini-0275:phyxminiproblems-problemphyxmini0275-surfacecondition}
  \lean{PhyXMiniProblems.ProblemPhyXMini0275.SurfaceCondition}
  Contact condition for the horizontal surface
\end{definition}
\begin{proof}
  This is the declared model definition of Surface Condition.
\end{proof}
\begin{definition}[Motion Regime]
  \label{def:physics:phyx-mini-0275:phyxminiproblems-problemphyxmini0275-motionregime}
  \lean{PhyXMiniProblems.ProblemPhyXMini0275.MotionRegime}
  The motion model stated in the problem
\end{definition}
\begin{proof}
  This is the declared model definition of Motion Regime.
\end{proof}
\begin{definition}[Coordinate Direction]
  \label{def:physics:phyx-mini-0275:phyxminiproblems-problemphyxmini0275-coordinatedirection}
  \lean{PhyXMiniProblems.ProblemPhyXMini0275.CoordinateDirection}
  Direction of motion along the one-dimensional coordinate axis
\end{definition}
\begin{proof}
  This is the declared model definition of Coordinate Direction.
\end{proof}
\begin{definition}[Graph Axis]
  \label{def:physics:phyx-mini-0275:phyxminiproblems-problemphyxmini0275-graphaxis}
  \lean{PhyXMiniProblems.ProblemPhyXMini0275.GraphAxis}
  The two axes of the supplied graph
\end{definition}
\begin{proof}
  This is the declared model definition of Graph Axis.
\end{proof}
\begin{definition}[Axis Quantity]
  \label{def:physics:phyx-mini-0275:phyxminiproblems-problemphyxmini0275-axisquantity}
  \lean{PhyXMiniProblems.ProblemPhyXMini0275.AxisQuantity}
  Physical quantity carried by a graph axis
\end{definition}
\begin{proof}
  This is the declared model definition of Axis Quantity.
\end{proof}
\begin{definition}[Figure Label]
  \label{def:physics:phyx-mini-0275:phyxminiproblems-problemphyxmini0275-figurelabel}
  \lean{PhyXMiniProblems.ProblemPhyXMini0275.FigureLabel}
  Text and symbolic labels visible in the primary image
\end{definition}
\begin{proof}
  This is the declared model definition of Figure Label.
\end{proof}
\begin{definition}[Force Position Graph]
  \label{def:physics:phyx-mini-0275:phyxminiproblems-problemphyxmini0275-forcepositiongraph}
  \lean{PhyXMiniProblems.ProblemPhyXMini0275.ForcePositionGraph}
  \uses{def:physics:phyx-mini-0275:phyxminiproblems-problemphyxmini0275-graphaxis, def:physics:phyx-mini-0275:phyxminiproblems-problemphyxmini0275-axisquantity, def:physics:phyx-mini-0275:phyxminiproblems-problemphyxmini0275-figurelabel}
  The calibrated scalar readout of the supplied force--position graph. The function records the plotted signed force component in newtons at a position read in metres; it is a measurement projection, not the physical force type
\end{definition}
\begin{proof}
  This is the declared model definition of Force Position Graph.
\end{proof}
\begin{definition}[Block Spring Oscillator Setup]
  \label{def:physics:phyx-mini-0275:phyxminiproblems-problemphyxmini0275-blockspringoscillatorsetup}
  \lean{PhyXMiniProblems.ProblemPhyXMini0275.BlockSpringOscillatorSetup}
  \uses{def:physics:phyx-mini-0275:phyxminiproblems-problemphyxmini0275-massquantity, def:physics:phyx-mini-0275:phyxminiproblems-problemphyxmini0275-lengthquantity, def:physics:phyx-mini-0275:phyxminiproblems-problemphyxmini0275-timequantity, def:physics:phyx-mini-0275:phyxminiproblems-problemphyxmini0275-velocityquantity, def:physics:phyx-mini-0275:phyxminiproblems-problemphyxmini0275-forcecomponentquantity, def:physics:phyx-mini-0275:phyxminiproblems-problemphyxmini0275-springstiffnessquantity, def:physics:phyx-mini-0275:phyxminiproblems-problemphyxmini0275-restoringelement, def:physics:phyx-mini-0275:phyxminiproblems-problemphyxmini0275-surfacecondition, def:physics:phyx-mini-0275:phyxminiproblems-problemphyxmini0275-motionregime, def:physics:phyx-mini-0275:phyxminiproblems-problemphyxmini0275-coordinatedirection, def:physics:phyx-mini-0275:phyxminiproblems-problemphyxmini0275-forcepositiongraph}
  The dimensionful block--spring system, its trajectory observables, and the supplied calibrated graph. Physlib's scalar harmonic oscillator is retained as the underlying ideal model and is linked to the dimensionful mass and stiffness by the governing-law interface below
\end{definition}
\begin{proof}
  This is the declared model definition of Block Spring Oscillator Setup.
\end{proof}
\begin{definition}[Matches Problem Description]
  \label{def:physics:phyx-mini-0275:phyxminiproblems-problemphyxmini0275-matchesproblemdescription}
  \lean{PhyXMiniProblems.ProblemPhyXMini0275.MatchesProblemDescription}
  \uses{def:physics:phyx-mini-0275:phyxminiproblems-problemphyxmini0275-massinkilograms, def:physics:phyx-mini-0275:phyxminiproblems-problemphyxmini0275-lengthinmeters, def:physics:phyx-mini-0275:phyxminiproblems-problemphyxmini0275-timeinseconds, def:physics:phyx-mini-0275:phyxminiproblems-problemphyxmini0275-velocityinmeterspersecond, def:physics:phyx-mini-0275:phyxminiproblems-problemphyxmini0275-blockspringoscillatorsetup}
  Problem-statement data. In particular, no numerical initial speed or kinetic energy is supplied: only its positive direction is known
\end{definition}
\begin{proof}
  This is the declared model definition of Matches Problem Description.
\end{proof}
\begin{definition}[Matches Primary Figure]
  \label{def:physics:phyx-mini-0275:phyxminiproblems-problemphyxmini0275-matchesprimaryfigure}
  \lean{PhyXMiniProblems.ProblemPhyXMini0275.MatchesPrimaryFigure}
  \uses{def:physics:phyx-mini-0275:phyxminiproblems-problemphyxmini0275-lengthquantity, def:physics:phyx-mini-0275:phyxminiproblems-problemphyxmini0275-lengthinmeters, def:physics:phyx-mini-0275:phyxminiproblems-problemphyxmini0275-forcecomponentinnewtons, def:physics:phyx-mini-0275:phyxminiproblems-problemphyxmini0275-blockspringoscillatorsetup}
  Readouts taken from the primary image. The endpoint positions are treated as the turning positions of the depicted oscillation, so the positive endpoint is the physical amplitude. The plotted signed scalar curve is explicitly linked to the dimensionful net-force observable throughout the shown range
\end{definition}
\begin{proof}
  This is the declared model definition of Matches Primary Figure.
\end{proof}
\begin{definition}[Has Physical Parameters]
  \label{def:physics:phyx-mini-0275:phyxminiproblems-problemphyxmini0275-hasphysicalparameters}
  \lean{PhyXMiniProblems.ProblemPhyXMini0275.HasPhysicalParameters}
  \uses{def:physics:phyx-mini-0275:phyxminiproblems-problemphyxmini0275-massinkilograms, def:physics:phyx-mini-0275:phyxminiproblems-problemphyxmini0275-lengthinmeters, def:physics:phyx-mini-0275:phyxminiproblems-problemphyxmini0275-springstiffnessinnewtonspermeter, def:physics:phyx-mini-0275:phyxminiproblems-problemphyxmini0275-blockspringoscillatorsetup}
  Positivity conditions selecting the physical oscillator branch
\end{definition}
\begin{proof}
  This is the declared model definition of Has Physical Parameters.
\end{proof}
\begin{definition}[Satisfies Undamped Simple Harmonic Oscillator Laws]
  \label{def:physics:phyx-mini-0275:phyxminiproblems-problemphyxmini0275-satisfiesundampedsimpleharmonicoscillatorlaws}
  \lean{PhyXMiniProblems.ProblemPhyXMini0275.SatisfiesUndampedSimpleHarmonicOscillatorLaws}
  \uses{def:physics:phyx-mini-0275:phyxminiproblems-problemphyxmini0275-lengthquantity, def:physics:phyx-mini-0275:phyxminiproblems-problemphyxmini0275-timequantity, def:physics:phyx-mini-0275:phyxminiproblems-problemphyxmini0275-massinkilograms, def:physics:phyx-mini-0275:phyxminiproblems-problemphyxmini0275-lengthinmeters, def:physics:phyx-mini-0275:phyxminiproblems-problemphyxmini0275-velocityinmeterspersecond, def:physics:phyx-mini-0275:phyxminiproblems-problemphyxmini0275-forcecomponentinnewtons, def:physics:phyx-mini-0275:phyxminiproblems-problemphyxmini0275-springstiffnessinnewtonspermeter, def:physics:phyx-mini-0275:phyxminiproblems-problemphyxmini0275-energyinjoules, def:physics:phyx-mini-0275:phyxminiproblems-problemphyxmini0275-blockspringoscillatorsetup}
  Governing laws for an undamped one-dimensional simple harmonic oscillator. They are the dimensionful-readout counterparts of Physlib's linear force, kinetic-energy, potential-energy, and total-energy definitions. The final numerical maximum and answer choice do not occur in this interface
\end{definition}
\begin{proof}
  This is the declared model definition of Satisfies Undamped Simple Harmonic Oscillator Laws.
\end{proof}
\begin{definition}[Is Maximum Kinetic Energy]
  \label{def:physics:phyx-mini-0275:phyxminiproblems-problemphyxmini0275-ismaximumkineticenergy}
  \lean{PhyXMiniProblems.ProblemPhyXMini0275.IsMaximumKineticEnergy}
  \uses{def:physics:phyx-mini-0275:phyxminiproblems-problemphyxmini0275-timequantity, def:physics:phyx-mini-0275:phyxminiproblems-problemphyxmini0275-energyinjoules, def:physics:phyx-mini-0275:phyxminiproblems-problemphyxmini0275-blockspringoscillatorsetup}
  An energy is the attained maximum kinetic energy when it is the actual kinetic energy at some time and its joule readout is a greatest element of the range of kinetic-energy readouts along the motion
\end{definition}
\begin{proof}
  This is the declared model definition of Is Maximum Kinetic Energy.
\end{proof}
\begin{definition}[Answer Choice]
  \label{def:physics:phyx-mini-0275:phyxminiproblems-problemphyxmini0275-answerchoice}
  \lean{PhyXMiniProblems.ProblemPhyXMini0275.AnswerChoice}
  Labels of the four displayed multiple-choice answers
\end{definition}
\begin{proof}
  This is the declared model definition of Answer Choice.
\end{proof}
\begin{definition}[joules]
  \label{def:physics:phyx-mini-0275:phyxminiproblems-problemphyxmini0275-answerchoice-joules}
  \lean{PhyXMiniProblems.ProblemPhyXMini0275.AnswerChoice.joules}
  \uses{def:physics:phyx-mini-0275:phyxminiproblems-problemphyxmini0275-answerchoice}
  Joule value printed beside each answer choice
\end{definition}
\begin{proof}
  This is the declared model definition of joules.
\end{proof}
\begin{definition}[recorded Answer Choice]
  \label{def:physics:phyx-mini-0275:phyxminiproblems-problemphyxmini0275-recordedanswerchoice}
  \lean{PhyXMiniProblems.ProblemPhyXMini0275.recordedAnswerChoice}
  \uses{def:physics:phyx-mini-0275:phyxminiproblems-problemphyxmini0275-answerchoice}
  Dataset metadata recording choice C; this is not a theorem premise
\end{definition}
\begin{proof}
  This is the declared model definition of recorded Answer Choice.
\end{proof}
\begin{definition}[Rounds To Nearest Tenth]
  \label{def:physics:phyx-mini-0275:phyxminiproblems-problemphyxmini0275-roundstonearesttenth}
  \lean{PhyXMiniProblems.ProblemPhyXMini0275.RoundsToNearestTenth}
  \uses{def:physics:phyx-mini-0275:phyxminiproblems-problemphyxmini0275-energyinjoules}
  Conventional nearest-tenth rounding with half-ties rounded upward: a positive exact value in [displayed - 0.05, displayed + 0.05) prints as displayed
\end{definition}
\begin{proof}
  This is the declared model definition of Rounds To Nearest Tenth.
\end{proof}
\begin{definition}[Is Closest Displayed Answer]
  \label{def:physics:phyx-mini-0275:phyxminiproblems-problemphyxmini0275-isclosestdisplayedanswer}
  \lean{PhyXMiniProblems.ProblemPhyXMini0275.IsClosestDisplayedAnswer}
  \uses{def:physics:phyx-mini-0275:phyxminiproblems-problemphyxmini0275-energyinjoules, def:physics:phyx-mini-0275:phyxminiproblems-problemphyxmini0275-answerchoice}
  A displayed choice is no farther from the exact energy than any option
\end{definition}
\begin{proof}
  This is the declared model definition of Is Closest Displayed Answer.
\end{proof}
\begin{lemma}[spring Stiffness In Newtons Per Meter eq two Hundred Fifty]
  \label{lem:physics:phyx-mini-0275:phyxminiproblems-problemphyxmini0275-springstiffnessinnewtonspermeter-eq-twohundredfifty}
  \lean{PhyXMiniProblems.ProblemPhyXMini0275.springStiffnessInNewtonsPerMeter_eq_twoHundredFifty}
  \uses{def:physics:phyx-mini-0275:phyxminiproblems-problemphyxmini0275-springstiffnessinnewtonspermeter, def:physics:phyx-mini-0275:phyxminiproblems-problemphyxmini0275-blockspringoscillatorsetup, def:physics:phyx-mini-0275:phyxminiproblems-problemphyxmini0275-matchesprimaryfigure, def:physics:phyx-mini-0275:phyxminiproblems-problemphyxmini0275-satisfiesundampedsimpleharmonicoscillatorlaws}
  The endpoint force and displacement determine the spring stiffness: k = F\_s / A = 75 N / 0.30 m = 250 N/m
\end{lemma}
\begin{proof}
  The conclusion follows from the governing relations and figure readouts named in the statement of spring Stiffness In Newtons Per Meter eq two Hundred Fifty.
\end{proof}
\begin{lemma}[maximum Kinetic Energy exact]
  \label{lem:physics:phyx-mini-0275:phyxminiproblems-problemphyxmini0275-maximumkineticenergy-exact}
  \lean{PhyXMiniProblems.ProblemPhyXMini0275.maximumKineticEnergy_exact}
  \uses{def:physics:phyx-mini-0275:phyxminiproblems-problemphyxmini0275-lengthinmeters, def:physics:phyx-mini-0275:phyxminiproblems-problemphyxmini0275-springstiffnessinnewtonspermeter, def:physics:phyx-mini-0275:phyxminiproblems-problemphyxmini0275-energyinjoules, def:physics:phyx-mini-0275:phyxminiproblems-problemphyxmini0275-blockspringoscillatorsetup, def:physics:phyx-mini-0275:phyxminiproblems-problemphyxmini0275-matchesproblemdescription, def:physics:phyx-mini-0275:phyxminiproblems-problemphyxmini0275-matchesprimaryfigure, def:physics:phyx-mini-0275:phyxminiproblems-problemphyxmini0275-hasphysicalparameters, def:physics:phyx-mini-0275:phyxminiproblems-problemphyxmini0275-satisfiesundampedsimpleharmonicoscillatorlaws, def:physics:phyx-mini-0275:phyxminiproblems-problemphyxmini0275-ismaximumkineticenergy}
  Energy conservation transfers the spring energy at the amplitude turning point to kinetic energy at equilibrium. Nonnegative spring potential energy also shows this equilibrium kinetic energy is a genuine attained maximum. The force--displacement triangle has area (1/2) F\_s A = (1/2)(75)(0.30) = 45/4 J
\end{lemma}
\begin{proof}
  The conclusion follows from the governing relations and figure readouts named in the statement of maximum Kinetic Energy exact.
\end{proof}
% --- Archon declaration topology end ---
% --- Archon physics formalization source end ---
